\section*{Chapitre \ref{ch:g12:age-of-the-earth} --- La datation radiométrique et l'âge de la Terre}

\begin{solution}{exo:g12:age-of-the-earth:1}
(a) un \omterm{def:g12:age-of-the-earth:closed-system}{système clos}~; (b) une \omterm{def:g12:age-of-the-earth:isochron}{isochrone}~; (c) l'\omterm{def:g12:age-of-the-earth:age}{âge de la Terre}.
\end{solution}

\begin{solution}{exo:g12:age-of-the-earth:2}
Plagioclase (0.35), hornblende (1.20), roche totale (3.10), feldspath
potassique (6.40), biotite (24.0). La biotite se place le plus haut~: c'est
elle qui contient le plus de rubidium~87 par rapport à son strontium~86~;
son rubidium a donc produit le plus de strontium~87 depuis la formation du
granite.
\end{solution}

\begin{solution}{exo:g12:age-of-the-earth:3}
$T = \ln 2/\lambda = 0.6931/\num{1.3972e-11} \approx \num{4.96e10}$ ans,
soit 49.6~milliards d'années~: $49.6/4.54 \approx 11$ fois l'\omterm{def:g12:age-of-the-earth:age}{âge de la
Terre}.
\end{solution}

\begin{solution}{exo:g12:age-of-the-earth:4}
$\lambda t = \num{1.3972e-11} \times \num{3.10e8} \approx 0.004331$, et
$e^{0.004331} - 1 \approx 0.004341$. Pour 480 millions d'années, $\lambda t
\approx 0.006707$ et $e^{\lambda t} - 1 \approx 0.006729$. Les deux valeurs
sont proches de $\lambda t$ lui-même, puisque $\lambda t$ est petit.
\end{solution}

\begin{solution}{exo:g12:age-of-the-earth:5}
Phillips (1860), puis Kelvin (1862), Joly (1899), Holmes (1927) et
Patterson (1956). $4540/100 \approx 45$ fois.
\end{solution}

\begin{solution}{exo:g12:age-of-the-earth:6}
$0.7080 + 10.0 \times 0.004341 \approx 0.7514$ aujourd'hui~; 0.7080 quand le
granite s'est formé, comme pour tous les minéraux du granite.
\end{solution}

\begin{solution}{exo:g12:age-of-the-earth:7}
(1) $m = (0.77525 - 0.71806)/(9.40 - 0.90) = 0.05719/8.50 \approx 0.006728$~;
$t = \ln(1.006728)/\num{1.3972e-11} \approx \num{4.80e8}$ ans, soit 480
millions d'années. (2) $0.71806 + 1.70 \times 0.006728 \approx 0.72950$ et $0.71806 + 4.20
\times 0.006728 \approx 0.74632$~: tous deux sont sur la droite. (3) Oui~:
le granite recoupe le schiste~; il est donc plus jeune, et 310 est
inférieur à 480 millions d'années.
\end{solution}

\begin{solution}{exo:g12:age-of-the-earth:8}
Tous ses points tomberaient au même endroit~: un seul point ne définit ni
droite ni pente. La biotite se trouve très à droite des autres~; une petite
erreur sur l'un des points modifie alors la pente de la droite qui la
relie au plagioclase bien moins qu'elle ne le ferait pour deux points
proches l'un de l'autre.
\end{solution}

\begin{solution}{exo:g12:age-of-the-earth:9}
(1) Son rubidium n'a pas changé~: elle garde son $x$~; elle a moins de
strontium~87~: son $y$ diminue. Le point descend à la verticale. (2) La
droite qui la relie au plagioclase est moins pentue, et donne un âge trop
jeune. (3) La biotite n'est plus sur la droite des autres minéraux~: les
points ne sont pas alignés, signe qu'un échantillon n'est pas resté un
\omterm{def:g12:age-of-the-earth:closed-system}{système clos}.
\end{solution}

\begin{solution}{exo:g12:age-of-the-earth:10}
(1) $4540/96 \approx 47$ fois. (2) Si du sel quitte la mer, par exemple
enfoui dans les sédiments, la mer contient moins de sel que les fleuves
n'en ont apporté au total~: diviser le sel présent par l'apport annuel
donne une durée trop courte.
\end{solution}

\begin{solution}{exo:g12:age-of-the-earth:11}
$m = (34.86 - 10.34)/(50.28 - 9.46) = 24.52/40.82 \approx 0.601$, entre
les pentes pour 4.55 (0.590) et 4.60~milliards d'années (0.610)~: environ
4.58~milliards d'années.
\end{solution}

\begin{solution}{exo:g12:age-of-the-earth:12}
(1) Le basalte, au rapport le plus faible, provient du manteau~; le
granite, au rapport le plus élevé, au moins en partie de la fusion de la
croûte. (2) Le rubidium passe dans le liquide quand le manteau fond~; il
s'est donc concentré dans la croûte, qui a produit davantage de
strontium~87 au fil des âges. (3) Non~: des minéraux qui cristallisent à
partir d'un même magma partagent son rapport initial, ce qui n'est pas le
cas de ces deux roches.
\end{solution}

\begin{solution}{exo:g12:age-of-the-earth:13}
(1) Rubidium~87~: $2^{-4.54/49.61} \approx 0.94$, il en reste environ
\qty{94}{\percent}~; uranium~238~: $2^{-4.54/4.47} \approx 0.49$, environ
la moitié. (2) $m =
e^{\lambda t} - 1 \approx \num{1.3972e-11} \times 10^6 \approx
\num{1.4e-5}$~: même la biotite du granite de Varden ($x = 24$) ne
s'élèverait que de $0.0003$ au-dessus du rapport initial. Il s'est formé si
peu de strontium~87 que la pente se distingue mal de zéro~: une horloge
lente convient aux roches anciennes.
\end{solution}

\begin{solution}{exo:g12:age-of-the-earth:14}
(1) $4.438 + 0.6026 \times 19.0 \approx 15.89$~; $15.5$ est inférieur de $0.39$,
soit environ \qty{2.4}{\percent}~: le sédiment est proche de la droite.
(2) 0.6026 est compris entre 0.590 (4.55) et 0.610 (4.60)~: environ
4.58~milliards d'années, dans l'intervalle de Patterson, de 4.48 à
4.62~milliards d'années.
\end{solution}

\begin{solution}{exo:g12:age-of-the-earth:15}
(1) Le stock ne perd plus alors que $46.7 - 24.4 = \qty{22.3}{TW}$ net~; le
même stock dure donc $46.7/22.3 \approx 2.1$ fois plus longtemps. (2) $96 \times 2.1 \approx
200$ millions d'années, encore plus de vingt fois trop peu. (3) La
\omterm{def:g12:internal-heat:convection}{convection mantellique}, qui remonte la chaleur des profondeurs et maintient
le gradient près de la surface élevé bien plus longtemps que la conduction
seule.
\end{solution}

\begin{solution}{pb:g12:age-of-the-earth:1}
\textbf{1.} La biotite ($x = 24.0$).

\textbf{2.} Ils ont cristallisé à partir d'un même magma, et les minéraux
qui cristallisent à partir d'un même liquide incorporent son strontium avec
son rapport isotopique.

\textbf{3.} Le strontium~86 est stable et n'est produit par aucune
désintégration~: dans un \omterm{def:g12:age-of-the-earth:closed-system}{système clos}, son nombre ne change pas~; c'est
donc un étalon fixe.

\textbf{4.} Le strontium~87, produit par la désintégration du
rubidium~87~; le rubidium~87 a diminué et le strontium~86 est resté le
même.

\textbf{5.} Une droite horizontale d'équation $y = 0.7080$~: même rapport,
$x$ différents.

\textbf{6.} $(0.81218 - 0.70952)/(24.0 - 0.35) = 0.10266/23.65 \approx
0.004341$.

\textbf{7.} $(0.72146 - 0.70952)/(3.10 - 0.35) = 0.01194/2.75 \approx
0.004342$.

\textbf{8.} Elles concordent à 1 pour 4000 près. La première est la plus
fiable~: ses deux points sont très éloignés, si bien que l'arrondi des
rapports compte moins.

\textbf{9.} $0.70952 - 0.35 \times 0.004341 \approx 0.7080$.

\textbf{10.} $0.7080 + 6.40 \times 0.004341 \approx 0.73578$~: exactement
la valeur mesurée.

\textbf{11.} $e^{\lambda t} = 1 + m = 1.004341$.

\textbf{12.} $\lambda t = \ln 1.004341 \approx 0.004331$.

\textbf{13.} $t = 0.004331/\num{1.3972e-11} \approx \num{3.10e8}$ ans.

\textbf{14.} Oui~: le granite recoupe le schiste~; d'après le principe de
recoupement, il est donc plus jeune, et 310 est inférieur à 480 millions
d'années.

\textbf{15.} 310 millions d'années est compris entre 358.9 et 298.9
millions d'années~: c'est le Carbonifère.

\textbf{16.} Le rapport $^{87}\mathrm{Sr}/^{86}\mathrm{Sr}$ du magma au
moment où le granite a cristallisé.

\textbf{17.} Le granite~: $0.7080 > 0.7035$.

\textbf{18.} Le rubidium passe dans le liquide quand le manteau fond~; il
s'est donc concentré dans la croûte~; avec davantage de rubidium, la croûte
a produit davantage de strontium~87.

\textbf{19.} Que le magma provenait, au moins en partie, de la fusion de
la \omterm{def:g12:continental-crust:continental-crust}{croûte continentale}, et non du seul manteau.

\textbf{20.} Le granite de Varden a \textbf{310 millions d'années}.
\end{solution}
